\section{Details for Section~\ref{sec:FlexMDM_inference}}
\label{sec:appendix_theory_inference}

\subsection{Precise detail on the inference algorithms}
In this section, we provide details on the unmasking steps of vanilla and adaptive inference for FlexMDM, summarized in Algorithm~\ref{alg:inference}, Subroutine 2. Suppose at inference time we are given the discretization step size $\tau$, a partially observed sequence $X_{t_k}$, and the current time step $t_{k}$.

\paragraph{Vanilla inference.} 
For each masked position $i$ (i.e., $X_{t_k}^i = \mask$) and each clean token $v \in \Sigma$, we sample unmasking events from a Poisson distribution $\mathrm{Poisson}(R_v \tau)$, where $R_v$ is the unmasking rate toward token $v$. 
Concretely, $R_v = \frac{\dot{\beta}_{t_k}}{1-\beta_{t_k}} \cdot f_\theta(X_{t_k}, t_k)[i,v]$,
so that the event count is distributed as
$k_v \sim \mathrm{Poi}\left(\tau \cdot \tfrac{\dot{\beta}_{t_k}}{1-\beta_{t_k}} \cdot f_\theta(X_{t_k}, t_k)[i,v]\right)$. A masked position is unmasked only if \emph{exactly one} token $v$ produces a count $k_v = 1$ while all others produce zero. This tau-leaping scheme batches all events that occur within the interval $[t_k, t_k+\tau]$.

\paragraph{Adaptive inference.} 
We first draw the number of tokens to unmask, denoted by an integer $K$. While Proposition~\ref{thm:inference} (Theorem~\ref{thm:anyorder}) shows that the choice of $K$ does not affect the theoretical guarantees, in practice, we set $K$ to match the expected number of unmasked tokens under vanilla inference yields stable behavior. Accordingly, we sample
$K \sim \mathrm{Poi}\!\left(\tau \cdot \tfrac{\dot{\beta}_{t_k}}{1-\beta_{t_k}} \cdot \#\{\text{masked tokens in } X_{t_k}\}\right)$.


Next, we compute a confidence score for each masked position, based on heuristics such as:
\begin{itemize}[leftmargin=*,itemsep=0pt,topsep=0pt]
    \item \underline{Top-K probability}~\citep{chang2022maskgit,zheng2024masked}: For state $x$ at time $t$, the confidence at position $i$ is given by $\max_{v \in \Sigma} f_\theta(x, t)[i,v]$.
    \item \underline{Top-K probability with sliding window}: We further restrict sampling to the leftmost $L$ tokens, where
    \[
    L = \min(\lfloor \gamma_1 \cdot L \rfloor, \gamma_2),
    \]
    with $\gamma_1$ and $\gamma_2$ hyperparameters. This approach is related to semi-autoregressive strategies used in \citet{nie2025large}.
\end{itemize}

Finally, we select the subset of positions to unmask as the Top-K masked indices with the highest confidence scores.


\subsection{Proof of FlexMDM's any-order inference capability}
\subsubsection{Proof preliminaries} \label{sec:appendix_unmasking_posterior}
\paragraph{Form of posterior.}
We first compute, for each $x_t^i$, the probabilities of being masked or deleted in equation~\eqref{eqn::FlexMDM_interpolant}. This follows from a straightforward calculation using the joint distribution of $(T_1^i,T_2^i)$:
\begin{align*}
    p(x_t^i=\text{(empty)}) &= p(T_1^i>t) = 1-\alpha_t, \\
    p(x_t^i=\mask) &= p(T_1^i \le t, T_2^i>t) =\int_{t}^1 \int_{0}^t\left( \frac{\dot{\beta}_s}{1-\beta_u} \times \dot{\alpha}_s\right) ds du =\colon 1-\gamma_t  .
\end{align*}
Here we define $\gamma_t$ as $1-p(x_t^i=\mask)$. Therefore, the process in equation~\eqref{eqn::FlexMDM_interpolant} is equivalent to observing a partially masked subsequence $x_t$ obtained by sampling $x_1 \sim p$ and, for each position of $x_1$, independently deleting it with probability $1-\alpha_t$, masking it with probability $1-\gamma_t$, or leaving it unchanged with probability $\alpha_t+\gamma_t-1$. Note that $\alpha_t$ and $\gamma_t$ both increase from $0$ to $1$ as $t$ increases from $0$ to $1$.

The posterior is given by
\begin{align*}
    \MoveEqLeft p(x_1 = x^*\mid x_t = x) \\
    &\propto p(x^*) \cdot p(x_t = x \mid x_1 = x^*) \\
    &= p(x^*) \cdot (1 - \alpha_t)^{\length{x^*} - \length{x}} (1 - \gamma_t)^{\#\mathrm{mask}(x)} (\alpha_t + \gamma_t - 1)^{\#\mathrm{unmask}(x)} \cdot \#\{s: x \subseteq x^*|_s\} \\
    &\propto p(x^*) \cdot (1 - \alpha_t)^{\length{x^*}-\length{x}} \cdot \#\{s: x\subseteq x^*|_s\}\,. \label{eq:posterior}
\end{align*}

Importantly, as in the vanilla MDM setting, the posterior does not depend on the unmasking schedule $(\gamma_t)$ (thus $\beta_t$), which will enable us to perform unmasking in adaptively chosen positions. Note also that if all sequences in the support of $p$ were of the same length, this posterior would also be independent of $(\alpha_t)$; while we do not prove it, in this case this would allow us to choose an arbitrary order of unmaskings and insertions.

\paragraph{Extension of posterior to $t = 1$.}

Motivated by the form of the posterior above, we define the following:
\begin{equation}
    q_t(x^* \mid x) \propto \begin{cases}
        p(x^*) \cdot \mathbf{1}_{x\subseteq x^*} & \text{if} \ t = 1 \\
        p(x^*) \cdot (1 - \alpha_t)^{\length{x^*} - \length{x}}\cdot \#\{s: x\subseteq x^*|_s\} & \text{otherwise} 
    \end{cases}
\end{equation}
Note that for $t < 1$, this is the same as $p(x_1 = x \mid x_t = x)$. We will denote the marginals of $q_t(\cdot \mid x)$ by $q^i_t(\cdot \mid x)$ for $v\in \Sigma$. The reason for extending the definition of the posterior to $t = 1$ is that in an adaptive FlexMDM sampler (see Definition~\ref{def:adaptive}), because we are entirely decoupling unmasking from the schedule of insertions, after the final insertion step the time parameter $t$ may be $1$ even though there are still tokens left to unmask. We will assume oracle access to $q_1(\cdot \mid x)$ as in practice these are simply the any-order marginals for $p$, and furthermore in practice these are already well-approximated by the learned posterior marginals $p(x^i_1\cdot \mid x_{1 - \delta} = x)$ for arbitrarily small $\delta > 0$.


\paragraph{Index-tracking variable.} Recall that in the definition of the joint interpolant we defined an index-tracking variable $s_t$ which essentially tracked which indices of $x_1$ correspond to the tokens in $x_t$. While our analysis below will not use the language of stochastic interpolants, we will still use the idea of tracking $s_t$, with slightly modified notation. Specifically, for any $0 \le t < 1$, we will use the notation $\Pr_{(x_1, s) \mid x_t = x}$ and $\E_{(x_1, s) \mid x_t = x}$ to denote probability and expectation with respect to the distribution given by sampling $x_1$ from $q_t(\cdot \mid x)$, and then sampling $s$ uniformly random from subsets for which $x\subseteq (x_1)|_s$. When we only care about the marginal distribution over $s$, we write $\Pr_{s\mid x_t = x}$ and $\E_{s\mid x_t = x}$. Given such a subset $s$ and $i\in\{0,\ldots,|s|-1\}$, we use $s_i$ to denote its $i$-th element in sorted order; as before, we also define the boundary values $s_{-1} = -1$ and $s_{\mathrm{len}(s)} = \mathrm{len}(x_1)$. The insertion expectations which we had denoted by $\E[s_t[i] - s_t[i-1] - 1]$ in the main body are thus given by $\E_{s\mid x_t = x}[s_i - s_{i-1} - 1]$ in the notation of this section.

\subsection{Formal guarantee for adaptive inference}

\begin{appdefinition}\label{def:adaptive}
    Given query access to the posterior marginals $q^i_t(\cdot \mid x_t = x)$ and to the insertion expectations $\E_{s\mid x_t = x}[s_i - s_{i-1} -1]$, an \emph{adaptive FlexMDM sampler} is any algorithm which produces a sequence of iterates $\hat{x}_{t_1},\ldots,\hat{x}_{t_n}$, where $0 = t_1 < \cdots < t_n = 1$, by starting at $\hat{x}_{t_1} = \varepsilon$ and arbitrarily alternating between steps of the following form:
    \begin{itemize}[leftmargin=*]
        \item \underline{\emph{Any-order} unmasking step}: Starting from $\hat{x}_{t_j}$, if $\mathrm{mask}(\hat{x}_{t_j})$ is nonempty, pick an arbitrary index $i$ therein (possibly probabilistically), sample $v$ from $q^i_{t_j}(\cdot \mid  \hat{x}_{t_j})$, and set $\hat{x}_{t_j} \gets \hat{x}_{t_j}[\hat{x}_{t_j}^i \gets v]$.
        \item \underline{Insertion step}: Starting from $\hat{x}_{t_j}$, run the CTMC with rate matrix
        \begin{equation}
            R^{\rm ins}_t(x,y) = \begin{cases} \E_{s\mid x_t = x}[s_i - s_{i-1} - 1]\cdot \frac{\dot{\alpha}_t}{1 - \alpha_t} & \text{if} \ y = \insertat{x}{i}{\mask} \\
            -\sum^{\length{x}}_{i=0} R^{\rm ins}_t(x,\insertat{x}{i}{\mask}) & \text{if} \ y = x \\
            0 & \text{otherwise}
            \end{cases} \label{eq:insrate}
        \end{equation} for $t_j \le t \le t_{j+1}$ to obtain $\hat{x}_{t_{j+1}}$. If $t_{j+1} = 1$, then apply any-order unmasking until $\mathrm{mask}(\hat{x}_{t_{j+1}})$ is empty, and terminate.
    \end{itemize}
\end{appdefinition}

Note that the rate matrix in the second bullet point above is identical to the one in the main body except that we only consider transitions given by insertions.

Formally, we will show the following:

\begin{appthm}\label{thm:anyorder}
    Any adaptive FlexMDM sampler for $p$ will generate a sequence of iterates $\hat{x}_{t_1},\ldots,\hat{x}_{t_n}$ such that $\hat{x}_{t_n}$ is exactly a sample from $p$.
\end{appthm}

\subsection{Proof of Theorem~\ref{thm:anyorder}}

To show that adaptive sampling works, we inductively prove an even stronger statement: at any intermediate step of the sampler after it has produced $\hat{x}_{t_j}$, the final output $\hat{x}_{t_n}$ is a sample from $q_{t_j}(\cdot \mid \hat{x}_{t_j})$.

The following two lemmas provide the inductive steps for unmasking and insertion respectively:

\begin{applem}[Inductive step for unmasking]\label{lem:unmasking_step}
    Let $0 \le t \le 1$ and let $x$ be a partially unmasked subsequence of length $m$. Let $\pi = (\pi_i)_{i\in \mathrm{mask}(x)}$ denote any distribution over masked indices of $x$.
    Suppose that one runs the following:
    \begin{enumerate}[leftmargin=*,itemsep=0pt,topsep=0pt]
        \item Sample index $i$ from $\pi$ 
        \item Sample $v$ from the posterior marginal $q^i_t(\cdot \mid x_t = x)$
        \item Sample from $q_t(\cdot \mid x[x^i\gets v])$.
    \end{enumerate}
    The output of this procedure is a sample $x^*$ from $q_t(\cdot \mid x)$.
\end{applem}

\begin{applem}[Inductive step for insertion]\label{lem:insertion_step}
    Let $0 \le t < 1$ and let $x$ be a partially unmasked subsequence of length $m$. Let $0\le h \le 1 - t$ be any duration of time. Suppose that one runs the following:
    \begin{enumerate}[leftmargin=*,itemsep=0pt,topsep=0pt] 
        \item Starting from $x$, run the CTMC with rate matrix
        given by Eq.~\eqref{eq:insrate} for time $h$ to obtain $x'$
        \item Sample from $q_{t+h}(\cdot \mid x')$.
    \end{enumerate}
    The output of this procedure is a sample from the posterior $q_t(\cdot \mid x)$.
\end{applem}

\noindent We defer the proofs of these to Sections~\ref{sec:unmasking_step} and~\ref{sec:insertion_step} below. The idea for the former is identical to the proof of the folklore fact that vanilla MDMs can sample in any order~\citep{kim2025train}. The proof for the latter is more involved and involves explicitly verifying that the Kolmogorov \emph{backward} equation is satisfied by the rate matrix we have constructed.

Here we verify that these Lemmas are enough to establish Theorem~\ref{thm:anyorder}.

\begin{proof}[Proof of Theorem~\ref{thm:anyorder}]
    We show more generally that starting from any intermediate time step $\hat{x}_{t_j}$ (not just $j = 1$), any adaptive FlexMDM sampler outputs a sample from $q_{t_j}(\cdot \mid \hat{x}_{t_j})$. We do this by inducting on the total number of insertion steps that remain.
    
    As the base case for the induction, if no more insertion steps remain, then we must have $t_j = 1$. In this case, we can further induct on the number of unmasking steps and apply Lemma~\ref{lem:unmasking_step} with $t$ therein set to $1$ to conclude that the final output is a sample from $q_1(\cdot \mid \hat{x}_{t_j})$.

    For the inductive step, we have $t_j < 1$ and suppose we have shown that for any FlexMDM sampler that makes at most $m$ insertion steps, starting from any $\hat{x}_{t_j}$ at intermediate time $t_j$, it samples from $q_{t_j}(\cdot\mid \hat{x}_{t_j})$. Now consider a FlexMDM sampler that makes at most $m + 1$ insertion steps starting from $\hat{x}_{t_j}$ at intermediate time $t_j$. If in the next step it performs an insertion step, i.e. it runs the CTMC with rate matrix defined above for total time $h = t_{j+1} - t_j$, then by Lemma~\ref{lem:insertion_step} and the inductive hypothesis, it samples from $q_{t_j}(\cdot \mid \hat{x}_{t_j})$. Alternatively, suppose the sampler performs some sequence of $\ell$ unmasking steps before performing an insertion step. Then by further inducting on $\ell$, we conclude by Lemma~\ref{lem:unmasking_step} that the sampler eventually outputs a sample from $q_{t_j}(\cdot \mid \hat{x}_{t_j})$.

    Finally, the theorem follows from the special case where $t_j = 0$ and $\hat{x}_{t_j} = \varepsilon$.
\end{proof}


\subsection{Proof of Lemma~\ref{lem:unmasking_step}}
\label{sec:unmasking_step}


\begin{proof}
    Fix any index $i \in \mathrm{mask}(x)$. The marginal $q^i_t(\cdot\mid x)$ is given by
    \begin{equation}
    \Pr_{(x_1,S)\mid x_t=x}[(x_1)|_{s_i} = v] = \frac{\sum_{x_1, S: (x_1)|_{s_i} = v} p(x) \cdot (1 - \alpha_t)^{\length{x_1} - \length{x}}\cdot \#\{S: x\subseteq (x_1)|_S\}}{\sum_{x_1} p(x_1) (1 - \alpha_t)^{\length{x_1} - \length{x}} \cdot \#\{S: x\subseteq (x_1)|_S\}}\,. \label{eq:marg}
    \end{equation}
    The posterior $q_t(\cdot \mid x[x^i\gets v])$ is given by
    \begin{equation}
        q_t(x^*\mid x[x^i\leftarrow v]) = \frac{p(x^*) \cdot (1 - \alpha_t)^{\length{x^*} - \length{x}}\cdot \#\{S: x[x^i\leftarrow v]\subseteq x^*|_S\}}{\sum_{x_1} p(x_1) \cdot (1 - \alpha_t)^{\length{x_1} - \length{x}}\cdot \#\{S: x[x^i\leftarrow v]\subseteq (x_1)|_S\}}\,. \label{eq:posteriorunmask}
    \end{equation}
    Note that the numerator of Eq.~\eqref{eq:marg} and the denominator of Eq.~\eqref{eq:posteriorunmask} are the same. So conditioned on unmasking index $i$, the above procedure outputs $x^*$ with probability
    \begin{align*}
        \MoveEqLeft \sum_z \Pr_{(x_1,S)\mid x_t = x}[(x_1)|_{s_i} = v]\cdot q_t(x^* \mid x[x^i\leftarrow v]) \\
        &= \frac{p(x^*) \cdot (1 - \alpha_t)^{\length{x^*} - \length{x}} \cdot \sum_z \#\{S: x[x^i\gets v]\subseteq x^*|_S\}}{\sum_{x_1} p(x_1) (1 - \alpha_t)^{\length{x_1} - \length{x}} \cdot \#\{S: x\subseteq (x_1)|_S\}} \\
        &= q_t(x^*\mid x)\,.
    \end{align*}
    This holds conditioned on unmasking any $i \in\mathrm{mask}(x)$, so regardless of the choice of $\pi$ over such positions, we will sample from the correct distribution $q_t(\cdot\mid x)$.
\end{proof}


\subsection{Proof of Lemma~\ref{lem:insertion_step}}
\label{sec:insertion_step}

\begin{proof}
    It suffices to show that the rate matrix satisfies the Kolmogorov \emph{backward} equation
    \begin{equation}
        \partial_t q_t(x^*\mid x) = -\sum^{\length{x}}_{i=0} R^{\rm ins}_t(x,\insertat{x}{i}{\mask}) q_t(x^*\mid \insertat{x}{i}{\mask}) - R^{\rm ins}_t(x,x) q_t(x^*\mid x)\,. \label{eq:kbe}
    \end{equation}
    First note that the rate $R^{\rm ins}_t(x, \insertat{x}{i}{\mask})$ can be expressed as 
    \begin{equation}
        \frac{\sum_{x_1} p(x_1) \cdot (1 - \alpha_t)^{\length{x_1} - \length{x}}\cdot \sum_{S: x\subseteq (x_1)|_S} (s_i - s_{i-1} - 1)}{\sum_{x_1} p(x_1) \cdot (1 - \alpha_t)^{\length{x_1} - \length{x}} \cdot \#\{S: x\subseteq (x_1)|_S\}}\cdot \frac{\dot{\alpha}_t}{1 - \alpha_t}\,.
    \end{equation}
    Furthermore, $\sum_i (s_i - s_{i-1} - 1) = \length{x_1} - \length{x}$, so
    \begin{multline}
        \sum_i R^{\rm ins}_t(x,\insertat{x}{i}{\mask}) \\
        = \frac{\sum_{x_1} p(x_1) \cdot (1 - \alpha_t)^{\length{x_1} - \length{x}}\cdot\#\{S: x\subseteq (x_1)|_S\}\cdot (\length{x_1} - \length{x})}{\sum_{x_1} p(x_1) \cdot (1 - \alpha_t)^{\length{x_1} - \length{x}} \cdot \#\{S: x\subseteq (x_1)|_S\}}\cdot \frac{\dot{\alpha}_t}{1 - \alpha_t}\,. \label{eq:sumrates}
    \end{multline}
    Let us compute $\partial_t q_t(x^*\mid x)$:
    \begin{align}
        \MoveEqLeft -\frac{p(x^*)\cdot (1 - \alpha_t)^{\length{x^*} - \length{x}}\cdot \#\{S: x\subseteq x^*|_S\}\cdot (\length{x^*}-\length{x})}{\sum_{x_1} p(x_1)\cdot (1 - \alpha_t)^{\length{x_1} - \length{x}}\cdot \#\{S: x\subseteq (x_1)|_S\}}\cdot \frac{\dot{\alpha}_t}{1 - \alpha_t} \nonumber \\ 
        \MoveEqLeft + \Bigl[\frac{\sum_{x_1} p(x_1) \cdot (1 - \alpha_t)^{\length{x_1}-\length{x}} \cdot \#\{S: x \subseteq (x_1)|_S\}\cdot (\length{x_1}-\length{x})}{\sum_{x_1} p(x_1) \cdot (1 - \alpha_t)^{\length{x_1}-\length{x}} \cdot \#\{S: x \subseteq (x_1)|_S\}}\cdot \frac{\dot{\alpha}_t}{1 - \alpha_t} \nonumber \\
        &\qquad\times \frac{p(x^*) \cdot (1 - \alpha_t)^{\length{x^*} - \length{x}} \cdot \#\{S: x \subseteq x^*|_S\}}{\sum_{x_1} p(x_1) \cdot (1 - \alpha_t)^{\length{x_1}-\length{x}} \cdot \#\{S: x \subseteq (x_1)|_S\}} \Bigr]\label{eq:partialtp}
    \end{align}
    Note that by Eq.~\eqref{eq:sumrates}, the term in the parentheses in Eq.~\eqref{eq:partialtp} is exactly 
    \begin{equation}
        \sum^{\length{x}}_{i=0} R^{\rm ins}_t(x, \insertat{x}{i}{\mask}) q_t(x^* \mid x) = -R^{\rm ins}_t(x, x) q_t(x^*\mid x)\,,
    \end{equation}
    It remains to verify that the first term in Eq.~\eqref{eq:partialtp} is equal to $-\sum_i R^{\rm ins}_t(x,\insertat{x}{i}{\mask}) q_t(x^* \mid \insertat{x}{i}{\mask})$. To that end, we must show that
    \begin{multline}
        \sum^{\length{x}}_{i=0} \frac{\sum_{x_1} p(x_1) (1 - \alpha_t)^{\length{x_1} - \length{x}} \cdot \sum_{S: x\subseteq (x_1)|_S} (s_i - s_{i-1} - 1)}{\sum_{x_1} p(x_1) (1 - \alpha_t)^{\length{x_1} - \length{x}} \cdot \#\{S: \insertat{x_t}{i}{\mask}\subseteq (x_1)|_S\}}\cdot \#\{S: \insertat{x}{i}{\mask}\subseteq x^*|_S\} \\
        =  \#\{S: x\subseteq (x^*)|_S\} \cdot (\length{x^*} - \length{x}) \label{eq:mainequality}
    \end{multline}
    The key combinatorial step is as follows. For any $x_1$ in the support of $p$, consider a subset $S$ for which $x \subseteq (x_1)|_S$. Note that for every such $S$, we can uniquely associate exactly $s_i - s_{i-1} - 1$ different subsets $S'$ of size $|S| + 1$ for which $\insertat{x}{i}{\mask} \subseteq (x_1)|_{S'}$. Therefore, $\sum_{S: x\subseteq (x_1)|_S} (s_i - s_{i-1} - 1) = \#\{S: \insertat{x}{i}{\mask}\subseteq (x_1)|_S\}$, and the left-hand side of Eq.~\eqref{eq:mainequality} thus becomes
    \begin{equation}
        \sum^{\length{x}}_{i=0} \#\{S: \insertat{x}{i}{\mask} \subseteq x^*|_S\} = \sum^{\length{x}}_{i=0} \sum_{S: x\subseteq x^*|_S} (s_i - s_{i-1} - 1) = \sum_{S: x\subseteq (x_1)|_S} (\length{x^*} - \length{x_t})\,,
    \end{equation}
    which completes the proof of Eq.~\eqref{eq:mainequality}.
\end{proof}